\documentclass[12pt]{article}

\usepackage[a4paper,margin=1in]{geometry}
\usepackage{amsmath,amssymb}
\usepackage{setspace}

\setstretch{1.2}

\title{CSE400 – Fundamentals of Probability in Computing\\
Lecture 9: Uniform, Exponential, Laplace, and Gamma Random Variables}

\author{Instructor: Dhaval Patel, PhD}
\date{Date: February 2, 2026}

\begin{document}

\maketitle

\section*{1. Types of Continuous Random Variables}

The lecture introduces several types of continuous random variables and their probability distributions:

Uniform Random Variable

Exponential Random Variable

Laplace Random Variable

Gamma Random Variable

Additionally, examples and applications are provided for each distribution.

\section*{2. Uniform Random Variable}

\textbf{Definition}

A random variable $X$ is uniformly distributed over the interval $[a,b]$ if its probability density function (PDF) is

\[
f_X(x)=
\begin{cases}
\frac{1}{b-a}, & a \le x < b \\
0, & \text{elsewhere}
\end{cases}
\]

The cumulative distribution function (CDF) is

\[
F_X(x)=
\begin{cases}
0, & x<a \\
\dfrac{x-a}{b-a}, & a \le x < b \\
1, & x\ge b
\end{cases}
\]

The corresponding PDF and CDF graphs are illustrated in the lecture slides (Figure 3.8).

\section*{3. Example 1 – Uniform Random Variable}

\textbf{Problem}

The phase of a sinusoid $\theta$ is uniformly distributed over $[0,2\pi]$.

Its PDF is

\[
f_{\Theta}(\theta)=
\begin{cases}
\dfrac{1}{2\pi}, & 0\le \theta <2\pi \\
0, & \text{otherwise}
\end{cases}
\]

Find:

$Pr(\theta>3\pi/4)$

$Pr(\theta<\pi \mid \theta>3\pi/4)$

$Pr(\cos(\theta)<1/2)$

\textbf{Solution}

(a)

\[
Pr(\theta>3\pi/4)
\]

For a uniform random variable on $[0,2\pi]$

\[
Pr(a<\theta<b)=\frac{b-a}{2\pi}
\]

Thus

\[
Pr(\theta>3\pi/4)=\frac{2\pi-3\pi/4}{2\pi}=\frac{5}{8}
\]

(b) Conditional Probability

\[
Pr(A|B)=\frac{Pr(A\cap B)}{Pr(B)}
\]

Let

\[
A=(\theta<\pi)
\]

\[
B=(\theta>3\pi/4)
\]

\[
Pr(3\pi/4<\theta<\pi)=\frac{\pi-3\pi/4}{2\pi}=\frac{1}{8}
\]

\[
Pr(\theta<\pi|\theta>3\pi/4)=\frac{1/8}{5/8}=\frac{1}{5}
\]

(c)

\[
Pr(\cos(\theta)<1/2)
\]

\[
\cos(\theta)=1/2 \Rightarrow \theta=\frac{\pi}{3},\frac{5\pi}{3}
\]

Thus

\[
\cos(\theta)<1/2 \quad \text{for} \quad \frac{\pi}{3}<\theta<\frac{5\pi}{3}
\]

Therefore

\[
Pr(\cos(\theta)<1/2)=\frac{5\pi/3-\pi/3}{2\pi}
=\frac{4\pi/3}{2\pi}
=\frac{2}{3}
\]

\textbf{Applications of Uniform Distribution}

Examples include:

Phase of a sinusoidal signal when all phases in $[0,2\pi]$ are equally likely.

Random number generation between 0 and 1 for simulations.

Arrival time of a user within a known time window assuming no time preference.

\section*{4. Exponential Random Variable}

\textbf{Definition}

For $b>0$, the exponential random variable has

PDF

\[
f_X(x)=\frac{1}{b}e^{-x/b}u(x)
\]

CDF

\[
F_X(x)=[1-e^{-x/b}]u(x)
\]

where $u(x)$ is the unit step function.

\section*{5. Example 2 – Exponential Random Variable}

\textbf{Problem}

Let $X$ be an exponential random variable with

\[
f_X(x)=e^{-x}u(x)
\]

Find:

$Pr(3X<5)$

$Pr(3X<y)$

Let $Y=3X$. Find $f_Y(y)$.

\textbf{Solution}

Given

\[
f_X(x)=e^{-x}u(x)
\]

\[
F_X(x)=[1-e^{-x}]u(x)
\]

(a)

\[
Pr(3X<5)=Pr(X<5/3)
\]

\[
=F_X(5/3)
\]

\[
=1-e^{-5/3}
\]

(b)

\[
Pr(3X<y)=Pr(X<y/3)
\]

\[
=F_X(y/3)
\]

\[
=[1-e^{-y/3}]u(y)
\]

(c)

Let $Y=3X$

\[
F_Y(y)=[1-e^{-y/3}]u(y)
\]

Differentiating

\[
f_Y(y)=\frac{d}{dy}F_Y(y)
\]

\[
f_Y(y)=\frac{1}{3}e^{-y/3}u(y)
\]

Thus $Y$ is exponential with parameter

\[
\lambda=\frac{1}{3}
\]

\textbf{Applications of Exponential Distribution}

Time until radioactive decay of a particle.

Duration of an idle period in a wireless communication channel.

Time until the next earthquake in simplified seismic models.

\section*{6. Laplace Random Variable}

\textbf{Definition}

PDF

\[
f_X(x)=\frac{1}{2b}\exp\left(-\frac{|x|}{b}\right)
\]

CDF

\[
F_X(x)=
\begin{cases}
\frac{1}{2}\exp(x/b), & x<0 \\
1-\frac{1}{2}\exp(-x/b), & x\ge0
\end{cases}
\]

Application

The Laplace random variable is used to model the amplitude of speech signals.

\section*{7. Example 3 – Laplace Random Variable}

\textbf{Problem}

\[
f_W(w)=ce^{-2|w|}
\]

Find:

$c$

$Pr(-1<W<2)$

$Pr(W>0|-1<W<2)$

\textbf{Solution}

(a)

\[
\int_{-\infty}^{\infty}ce^{-2|w|}dw=1
\]

\[
2c\int_{0}^{\infty}e^{-2w}dw=1
\]

\[
2c\left(\frac{1}{2}\right)=1
\]

\[
c=1
\]

(b)

\[
Pr(-1<W<2)=\int_{-1}^{0}e^{2w}dw+\int_{0}^{2}e^{-2w}dw
\]

\[
=\frac{1}{2}(1-e^{-2})+\frac{1}{2}(1-e^{-4})
\]

\[
=1-\frac{1}{2}(e^{-2}+e^{-4})
\]

(c)

\[
Pr(W>0|-1<W<2)=\frac{Pr(0<W<2)}{Pr(-1<W<2)}
\]

\[
Pr(0<W<2)=\int_{0}^{2}e^{-2w}dw
\]

\[
=\frac{1}{2}(1-e^{-4})
\]

Thus

\[
Pr(W>0|-1<W<2)=
\frac{1-e^{-4}}{2-(e^{-2}+e^{-4})}
\]

\section*{8. Gamma Random Variable}

PDF

\[
f_X(x)=
\frac{(x/b)^{c-1}e^{-x/b}}{b\Gamma(c)}u(x)
\]

CDF

\[
F_X(x)=
\frac{\gamma(c,x/b)}{\Gamma(c)}u(x)
\]

\textbf{Definitions}

Gamma Function

\[
\Gamma(a)=\int_0^\infty e^{-t}t^{a-1}dt
\]

Incomplete Gamma Function

\[
\gamma(a,\beta)=\int_0^\beta e^{-t}t^{a-1}dt
\]

\textbf{Special Cases}

$b=2, c=0.5$ → Chi-square random variable

$c=1$ → Exponential random variable

with

$k=c, \theta=b$

\section*{9. Example 4 – Gamma Random Variable}

Let

\[
X \sim Gamma(\alpha,\lambda)
\]

with PDF

\[
f_X(x)=
\frac{\lambda^\alpha}{\Gamma(\alpha)}x^{\alpha-1}e^{-\lambda x},\quad x\ge0
\]

(a) Expected Value

\[
E[X]=\int_0^\infty x f_X(x)dx
\]

\[
E[X]=\frac{\alpha}{\lambda}
\]

(b) Variance

\[
Var(X)=E[X^2]-(E[X])^2
\]

\[
Var(X)=\frac{\alpha}{\lambda^2}
\]

\section*{10. Properties of the Gamma Function}

Prove the following

\[
\Gamma(n)=(n-1)!
\]

for $n=1,2,3,...$

\[
\Gamma(x+1)=x\Gamma(x)
\]

\[
\Gamma(1/2)=\sqrt{\pi}
\]

\section*{11. Applications of the Gamma Distribution}

Total service time for multiple independent tasks where each task duration is exponential.

Received signal energy computed from multiple independent signal samples.

Waiting time until the $k$-th event in a Poisson process.

\section*{12. Problem Solving Example – Movie Theater}

Suppose arrival time depends on travel time $T_r$ with PDF $f(t)$ and CDF $F(t)$.

Let $d$ be the departure time.

Arrival time

\[
Arrival=d+T_r
\]

Cost structure

Early arrival

\[
cost=c(T-d-T_r)
\]

Late arrival

\[
cost=k(d+T_r-T)
\]

Expected cost

\[
E[C(d)]=
\int_0^{T-d}c(T-d-t)f(t)dt
+
\int_{T-d}^{\infty}k(t+d-T)f(t)dt
\]

Optimal Departure Time

\[
-cF(T-d)+k[1-F(T-d)]=0
\]

Thus

\[
F(T-d^*)=\frac{k}{c+k}
\]

\[
d^*=T-F^{-1}\left(\frac{k}{c+k}\right)
\]

Interpretation

Depart such that the probability of arriving early equals

\[
\frac{k}{c+k}
\]

\section*{13. CDF Analysis Example}

Given

\[
F_X(x)=
\begin{cases}
0, & x<0 \\
x^2, & 0\le x \le1 \\
1, & x>1
\end{cases}
\]

PDF

\[
f_X(x)=\frac{d}{dx}F_X(x)=2x,\quad 0\le x\le1
\]

Mean

\[
\mu_X=\int_0^1 x(2x)dx
\]

\[
\mu_X=\frac{2}{3}
\]

Variance

\[
\sigma_X^2=\int_0^1 x^2(2x)dx-\mu_X^2
\]

\[
=\frac{1}{2}-\frac{4}{9}
\]

\[
\sigma_X^2=\frac{1}{18}
\]

Skewness

\[
c_s=\frac{E[(X-\mu)^3]}{\sigma^3}=-\frac{2\sqrt{2}}{5}
\]

Kurtosis

\[
c_k=\frac{E[(X-\mu)^4]}{\sigma^4}=\frac{12}{5}
\]

Conclusion

Distribution is left-skewed

Platykurtic (lighter tails than Gaussian)

\end{document}
```
